%%% FIELD proof-scm-counterpair
Put
\[
 \pi _0=(1/2,1/3,1/6),\qquad \pi _1=(1/3,1/3,1/3),
 \qquad q=(1/4,1/2,3/4),
\]
and write
\[
 s^+=(5/8,1/4,5/8),\qquad s^-=(3/8,3/4,3/8),
 \qquad t=(0,1/2,1).
\]
All calculations below are exact rational calculations.

First we verify the structural and independence premises in
\(\mathrm{def{:}wbc\text{-}class}\).  In either model, each potential response is
a measurable function of the independent shock specified in
\(\mathrm{def{:}failure\text{-}artifact}\): \(A\) uses \(U_A\), the entire
response \((H_0,H_1)\) uses \(U_H\), all \(Z_{a,h}\) use \(U_Z\), all
\(W_h\) use \(U_W\), all \(M_{a,h}\) use \(U_M\), and all
\(Y_{a,h,m,w}\) use \(U_Y\).  The pair \((U_H,U_Y)\) may be regarded as
the exogenous block for the district \(\{H,Y\}\); its two coordinates happen
to be independent, which is allowed.  The displayed structural substitutions
therefore have precisely the parent sets allowed by
\(\mathcal G_{\mathrm{WBC}}\), and the five consistency equalities hold by
construction.

Moreover, \(A\), being a function of \(U_A\), is independent of every
variable in
\(\{Y_{a,H_a,m,W_{H_a}},H_a\}\).  The variable \(M_{a',h}\) is a
function of \(U_M\), whereas
\(\{Y_{a,H_a,m,W_{H_a}},H_a\}\) is a function of
\((U_H,U_W,U_Y)\); this proves the cross-world independence.  The same
block independence gives \(M_{a,h}\mathbin{\perp\!\!\!\perp}(H,A)\).
Conditional on \((H,A)\), the realized mediator is a function of \(U_M\),
whereas \(Y_{A,H,m,W_H}\) is a function of \((U_W,U_Y)\); hence the
sequential-outcome independence holds.  Conditional on \((H,A)\), \(Z\)
uses only the independent shock \(U_Z\), proving
\(Z\mathbin{\perp\!\!\!\perp}(Y,M)\mid(H,A)\).  Conditional on \(H\),
\(W\) uses only \(U_W\), independently of the shocks determining
\((A,Z,M)\), proving
\(W\mathbin{\perp\!\!\!\perp}(A,Z,M)\mid H\).  Thus every listed
sequential, cross-world, consistency, and proxy-exclusion premise is
satisfied.

The observable denominator probabilities are
\[
 P(A=0,Z=z,M=m)=\frac12\frac12\frac12=\frac18>0,
 \qquad
 P(A=1,Z=z)=\frac12\frac12=\frac14>0.
\]
Every coordinate of \(\pi _0,\pi _1,s^+,s^-\) lies strictly between zero
and one.  Consequently, for every \(h,m\),
\[
 P(H_0=h)>0,\quad P(M_{1,h}=m)>0,\quad
 P(H=h,M=m,A=0)=\frac14\pi _0(h)>0,\quad
 P(H=h,A=1)=\frac16>0.
\]
These equations verify both target-relevant latent transport-overlap
restrictions.

It remains to verify latent-valid observed bridge existence.  In both models
take
\[
 h_0(0,0)=h_0(1,0)=\frac12,\qquad
 h_0(0,1)=-\frac12,\qquad h_0(1,1)=\frac32.
\]
Since \(W_h\) is Bernoulli with mean \(q_h\),
\[
 E\{h_0(W_h,0)\}=\frac12,\qquad
 E\{h_0(W_h,1)\}=-\frac12+2q_h=t_h.
\]
These are exactly the treatment-zero conditional outcome means for mediator
values zero and one, respectively, so the first latent bridge equations hold.
For the second bridge, conditional independence of \(W_h\) and
\(M_{1,h}\) yields
\[
 L_h^{\pm}:=E\{h_0(W_h,M_{1,h})\}
   =(1-s_h^{\pm})\frac12+s_h^{\pm}t_h.
\]
Direct substitution gives
\[
 (L_0^+,L_1^+,L_2^+)=(3/16,1/2,13/16),\qquad
 (L_0^-,L_1^-,L_2^-)=(5/16,1/2,11/16).
\]
For the two proposed second-stage bridges,
\[
\begin{aligned}
 E\{h_1^+(W_h)\}
   &=-\frac18(1-q_h)+\frac98q_h=L_h^+,\\
 E\{h_1^-(W_h)\}
   &= \frac18(1-q_h)+\frac78q_h=L_h^-.
\end{aligned}
\]
Thus the displayed vectors \(x^+\) and \(x^-\) in
\(\mathrm{def{:}failure\text{-}artifact}\) belong to their respective
latent bridge sets.  They also solve the observable equations.  Indeed, for
either vector the first two distinct row calculations are
\[
 \frac7{12}\frac12+\frac5{12}\frac12=\frac12,
 \qquad
 \frac7{12}\left(-\frac12\right)
   +\frac5{12}\frac32=\frac13,
\]
and the treatment-one row is
\[
 \frac14\left(\frac12+\frac12-\frac12+\frac32\right)
 -\frac12\{h_1(0)+h_1(1)\}=\frac12-\frac12=0,
\]
because the two coordinates of either \(h_1\) sum to one.  Hence
\(B(P^{\dagger})x^{\pm}=b(P^{\dagger})\), proving the required nonempty
latent-valid observed bridge intersection.  We have now verified every member
property of \(\mathfrak M_{\mathrm{WBC}}^{+}\).

We next verify observational equivalence.  The two models are identical in
the treatment-zero arm.  In the treatment-one arm, \(Z\) and \(Y\) are
independent Bernoulli variables with mean \(1/2\).  For either sign,
\[
 \frac13\sum_h s_h^{\pm}=\frac12,\qquad
 \frac13\sum_h q_h=\frac12,\qquad
 \frac13\sum_h q_hs_h^{\pm}=\frac14.
\]
Because \(W\) and \(M\) are binary, these three moments imply
\(P(W=w,M=m\mid A=1)=1/4\) for every \((w,m)\).  Thus the entire observed
law agrees and is the displayed \(P^{\dagger}\).  Every treatment-one cell
is positive.  In the treatment-zero arm, \(Z\) and \(M\) have probability
\(1/2\) at each value and every \(q_h\) is interior; for \(M=0\), both
outcome values have probability \(1/2\), while for \(M=1\) the positive
\(h=1\) mixture component assigns probability \(1/2\) to either outcome.
Therefore every cell of \(P^{\dagger}\) is strictly positive.

Under \(A=0\), \(M\) is independent of \(H\), and hence
\[
 P(W=1\mid A=0,Z=z,M=m)=\sum_h\pi _0(h)q_h=\frac5{12}.
\]
The treatment-zero outcome means are \(1/2\) for \(m=0\) and
\(\sum_h\pi _0(h)t_h=1/3\) for \(m=1\).  Under \(A=1\), the preceding
calculation gives \(P(W=w,M=m\mid Z=z,A=1)=1/4\) and
\(P(W=w\mid Z=z,A=1)=1/2\).  Substitution into
\(\mathrm{def{:}coupled\text{-}bridge\text{-}matrix}\) gives exactly the
displayed \(B(P^{\dagger})\), \(b(P^{\dagger})\), and
\(\ell(P^{\dagger})\).  The first, third, and fifth rows of the displayed
matrix are linearly independent and every other row duplicates one of these,
so its rank is three.  The vector
\[
 v=(0,0,0,0,1,-1)^{\top}
\]
satisfies \(B(P^{\dagger})v=0\), but
\[
 \ell(P^{\dagger})^{\top}v=\frac7{12}-\frac5{12}=\frac16\ne0.
\]
Every vector in the row space annihilates every vector in the null space;
therefore \(\ell(P^{\dagger})\notin
\operatorname{row}(B(P^{\dagger}))\), and the row-space criterion fails.

Finally, conditioning the nested target first on \(H_0=h\) gives
\[
 \psi(\mathcal S^{\pm})
 =\sum_{h=0}^2\pi _0(h)
 \left\{(1-s_h^{\pm})\frac12+s_h^{\pm}t_h\right\}.
\]
Consequently,
\[
\begin{aligned}
 \psi(\mathcal S^{\oplus})
 &=\frac12\frac3{16}+\frac13\frac12+\frac16\frac{13}{16}
   =\frac{19}{48},\\
 \psi(\mathcal S^{\ominus})
 &=\frac12\frac5{16}+\frac13\frac12+\frac16\frac{11}{16}
   =\frac7{16}.
\end{aligned}
\]
Their difference is \(1/24\), completing every assertion.
%%% FIELD proof-proxy-measurable-rescue
Fix an observed law \(P\) satisfying the premise and a map
\(f:\mathcal X\to\mathcal X\).  Let \(\mathcal S\) be any member of
\(\mathfrak M_{\mathrm{WBC}}^{+}\) inducing \(P\).  Define a bridge pair by
\[
 h_0^f(w,m)=f(w)\quad(w,m\in\mathcal X),
 \qquad h_1^f(w)=f(w)\quad(w\in\mathcal X),
\]
and let \(x_f\) be its six-coordinate vector.

The two observable denominator assumptions make every conditional moment in
the coupled bridge system well defined.  Since
\(P(Y=f(W)\mid A=0)=1\), for each \((z,m)\),
\[
\begin{aligned}
 E_P[Y\mid Z=z,M=m,A=0]
 &=E_P[f(W)\mid Z=z,M=m,A=0]\\
 &=\sum_w h_0^f(w,m)P(W=w\mid Z=z,M=m,A=0).
\end{aligned}
\]
For each \(z\),
\[
\begin{aligned}
 &\sum_{w,m}h_0^f(w,m)P(W=w,M=m\mid Z=z,A=1)\\
 &\qquad=\sum_w f(w)P(W=w\mid Z=z,A=1)\\
 &\qquad=\sum_w h_1^f(w)P(W=w\mid Z=z,A=1).
\end{aligned}
\]
By the definition of \(B(P)\) and \(b(P)\), these six equations say
\(B(P)x_f=b(P)\), so \(x_f\in\mathcal H_{\mathrm{obs}}(P)\).

We also verify latent validity rather than infer it from the observed
equations.  If \((h,m)\) is target relevant in the first latent bridge
equation, the latent treatment-zero overlap assumption gives
\(P_{\mathcal S}(H=h,M=m,A=0)>0\).  The observed almost-sure equality
\(Y=f(W)\) on \(\{A=0\}\) therefore remains true after conditioning on
this positive-probability event, and
\[
\begin{aligned}
 E_{P_{\mathcal S}}[Y\mid H=h,M=m,A=0]
 &=E_{P_{\mathcal S}}[f(W)\mid H=h,M=m,A=0]\\
 &=\sum_w h_0^f(w,m)
 P_{\mathcal S}(W=w\mid H=h,M=m,A=0).
\end{aligned}
\]
For every target-relevant \(h\), the second latent bridge equation is the
identity
\[
\begin{aligned}
 &\sum_{w,m}h_0^f(w,m)
 P_{\mathcal S}(W=w,M=m\mid H=h,A=1)\\
 &\qquad=\sum_w f(w)P_{\mathcal S}(W=w\mid H=h,A=1)\\
 &\qquad=\sum_w h_1^f(w)P_{\mathcal S}(W=w\mid H=h,A=1),
\end{aligned}
\]
whose conditionals are defined by the latent treatment-one overlap
assumption.  Hence
\(x_f\in\mathcal H_{\mathrm{lat}}(\mathcal S)\).  Applying
\(\mathrm{lem{:}latent\text{-}bridge\text{-}identifies}\) to this explicit
intersection point gives
\[
 \psi(\mathcal S)=\ell(P)^{\top}x_f
 =\sum_w P(W=w\mid A=0)f(w)
 =E_P[f(W)\mid A=0].
\]
The right side depends only on \(P\), and no row-space premise was used, so
the equality holds for every class member in the observed-law fiber.

For the exact nonconstant witness, modify either member of
\(\mathrm{def{:}failure\text{-}artifact}\) by setting
\(Y_{0,h,m,w}=w\) and leave every other response unchanged.  The consistency,
sequential, cross-world, exclusion, denominator, and overlap arguments in the
counterpair proof are unchanged: in particular, the new treatment-zero
potential outcome uses only \(U_W\) and is still independent of the mediator
shock \(U_M\).  The bridge pair
\[
 h_0(w,m)=w,\qquad h_1(w)=w
\]
is both observed-valid and latent-valid by the preceding calculation, so the
modified models remain in \(\mathfrak M_{\mathrm{WBC}}^{+}\).  The models
remain observationally equivalent because their treatment-zero mechanisms
are modified identically, while their only difference is still the
treatment-one mediator kernel and the treatment-one outcome remains
Bernoulli with mean \(1/2\).  The law of \((A,Z,W,M)\), and hence both
\(B(P)\) and \(\ell(P)\), is unchanged.  Thus the same null vector
\((0,0,0,0,1,-1)^{\top}\) proves
\(\ell(P)\notin\operatorname{row}(B(P))\).  On the other hand,
\[
 P(Y=1\mid A=0)=P(W=1\mid A=0)
 =\sum_h\pi _0(h)q_h=\frac5{12}.
\]
The already-proved general part, with \(f(w)=w\), therefore fixes the target
at \(5/12\) for every class member inducing this law.  Taking \(f\equiv0\)
gives the stated deterministic-zero corollary and proves that row span is not
universally necessary.
%%% FIELD proof-canonical-full-law-fiber
For an observed law \(P\), denote its full causal fiber by
\[
 \mathcal C(P)=\{\mathcal S\in\mathfrak M_{\mathrm{WBC}}^{+}:
 P_{\mathcal S}^O=P\}.
\]
We prove (17) by constructing mutually inverse representations at the level
of induced full-law target values.

\emph{From a class member to a feasible canonical point.}
Fix \(\mathcal S\in\mathcal C(P)\).  Because all endogenous alphabets are
finite, every realization of an exogenous shock determines a deterministic
response table.  Let \(R_A\) be the resulting treatment response, let
\((R_H,R_Y)\) be the joint response table for the district \(\{H,Y\}\),
and let \(R_Z,R_W,R_M\) be the other three response tables.  Compatibility
with \(\mathcal G_{\mathrm{WBC}}\) means that exogenous blocks belonging to
distinct districts are independent, while arbitrary dependence is allowed
inside the \(\{H,Y\}\) block.  Therefore their pushforward laws have the
form
\[
 \theta_A\otimes\theta_{HY}\otimes\theta_Z\otimes\theta_W
 \otimes\theta_M,
\]
on exactly the finite response-type spaces displayed in the theorem.  In
particular, the joint mass of a response tuple \(r\) is \(q_\theta(r)\).

Consistency and recursive structural substitution give
\[
 A=r_A,\quad H=r_H(r_A),\quad Z=r_Z(A,H),\quad W=r_W(H),\quad
 M=r_M(A,H),\quad Y=r_Y(A,H,M,W).
\]
These are precisely the recursions defining
\((a_r,h_r,z_r,w_r,m_r,y_r)\).  It follows cell by cell that
\[
 p_\theta(a,z,w,m,y)=P_{\mathcal S}^O(a,z,w,m,y)=P(a,z,w,m,y).
 \tag{22}
\]
The two denominator inequalities in (13) follow from the corresponding
member assumptions.

Independence of \(R_M\) from the other response blocks and consistency give
\[
\begin{aligned}
 u_h&=P_{\mathcal S}(H_0=h),&
 v_{hm}&=P_{\mathcal S}(M_{1,h}=m),\\
 j^0_{hm}&=P_{\mathcal S}(A=0,H=h,M=m),&
 j^1_h&=P_{\mathcal S}(A=1,H=h).
\end{aligned}
 \tag{23}
\]
Thus the two latent transport-overlap member assumptions imply (14).
Latent-valid observed bridge existence supplies an
\[
 x\in\mathcal H_{\mathrm{lat}}(\mathcal S)
      \cap\mathcal H_{\mathrm{obs}}(P).
\]
Observed validity is exactly \(B(P)x=b(P)\).  For every target-relevant
\((h,m)\), multiply the first conditional latent bridge equation by the
positive number \(j^0_{hm}\) in (23).  The result is (15), because
\(C^0_{hm}=\{A=0,H=h,M=m\}\) under the response recursion.  Similarly,
multiplying the second conditional latent bridge equation by the positive
number \(j^1_h\) gives (16).  Hence \((\theta,x)\in\mathcal F(P)\).

Finally, the same response table \(r_H(0)\) is used in both occurrences of
the hidden witness in
\[
 r_Y\!\left(0,r_H(0),r_M(1,r_H(0)),r_W(r_H(0))\right).
\]
Taking its expectation under \(q_\theta\) therefore gives, by the definition
of the nested potential outcome,
\[
 T(\theta)=E_{\mathcal S}
 [Y_{0,H_0,M_{1,H_0},W_{H_0}}]=\psi(\mathcal S).
 \tag{24}
\]
Equations (22)--(24) prove the inclusion of the left side of (17) in
\(\Gamma(P)\).

\emph{From a feasible canonical point to a class member.}
Conversely, fix \((\theta,x)\in\mathcal F(P)\).  On the finite product
response space, draw
\[
 R_A,\ (R_H,R_Y),\ R_Z,\ R_W,\ R_M
\]
independently across the five displayed blocks with marginal laws
\(\theta_A,\theta_{HY},\theta_Z,\theta_W,\theta_M\), respectively.  Define
the recursive SCM
\[
 A=R_A,\quad H=R_H(A),\quad Z=R_Z(A,H),\quad W=R_W(H),\quad
 M=R_M(A,H),\quad Y=R_Y(A,H,M,W).
 \tag{25}
\]
Only \(R_H\) and \(R_Y\) may be dependent, exactly realizing the district
\(\{H,Y\}\); all directed arguments in (25) are allowed parents in
\(\mathcal G_{\mathrm{WBC}}\).  Thus the constructed SCM, denoted
\(\mathcal S_\theta\), is compatible with that graph, and all five
consistency assumptions hold identically.

For completeness, we verify rather than assume the remaining structural
member properties.  The random variable \(A\) is a function of \(R_A\),
whereas \(\{Y_{a,H_a,m,W_{H_a}},H_a\}\) is a function of
\((R_H,R_Y,R_W)\), proving sequential treatment independence.  Each
\(M_{a,h}=R_M(a,h)\) is independent of \((H,A)\), and it is also independent
of \(\{Y_{a,H_a,m,W_{H_a}},H_a\}\), proving mediator independence and the
cross-world condition.  Conditional on \((H,A)\), the event conditioned on
depends only on \((R_H,R_A)\); the realized \(M\) then uses the independent
block \(R_M\), while \(Y_{A,H,m,W_H}\) uses \((R_Y,R_W)\).  This proves
the sequential-outcome condition.  In the same way, conditional on
\((H,A)\), \(Z\) uses the independent block \(R_Z\), so it is independent
of \((Y,M)\).  Conditional on \(H\), \(W\) uses the independent block
\(R_W\), whereas \((A,Z,M)\) uses \((R_A,R_H,R_Z,R_M)\); hence the two
proxy-exclusion conditions hold.  These conditional factorizations remain
valid on every conditioning event of positive probability, which is the
finite-law meaning of the asserted conditional independences.

By construction, the observed cell probabilities of (25) are
\(p_\theta(o)\), so the first feasibility constraint yields
\(P_{\mathcal S_\theta}^O=P\).  The strict inequalities in (13) give the two
observable denominator assumptions.  The equalities in (23) also hold for
the constructed SCM.  Therefore (14) gives both target-relevant latent
transport-overlap assumptions.

It remains only to prove bridge-supported membership.  Feasibility already
gives \(B(P)x=b(P)\), hence
\(x\in\mathcal H_{\mathrm{obs}}(P)\).  If \(u_hv_{hm}>0\), (14) gives
\(j^0_{hm}>0\).  Dividing (15) by this positive number converts its left side
to
\(E_{P_{\mathcal S_\theta}}[Y\mid H=h,M=m,A=0]\) and its right side to
\[
 \sum_w h_0(w,m)
 P_{\mathcal S_\theta}(W=w\mid H=h,M=m,A=0).
\]
This is exactly the first latent bridge equation.  If \(u_h>0\), (14) gives
\(j^1_h>0\); dividing (16) by \(j^1_h\) gives exactly the second latent
bridge equation.  Hence
\[
 x\in\mathcal H_{\mathrm{lat}}(\mathcal S_\theta)
      \cap\mathcal H_{\mathrm{obs}}(P),
\]
so the latent-valid bridge-existence assumption holds.  We have proved
\(\mathcal S_\theta\in\mathcal C(P)\).  Direct substitution in (25), using
one and the same coordinate \(R_H(0)\) throughout the nested counterfactual,
again gives
\[
 \psi(\mathcal S_\theta)=T(\theta).
 \tag{26}
\]
This proves the reverse inclusion in (17), and hence equality.

We now derive each characterization.  If \(P\) is induced by the class, its
fiber and therefore \(\Gamma(P)\) are nonempty.  Also
\(T(\theta)\in[0,1]\), since it is the expectation of a binary response.
For any nonempty subset \(G\subseteq\mathbb R\), every element of \(G\) is
equal if and only if \(\inf G=\sup G\).  Applying this elementary fact to
the exact equality (17) proves that the target is constant on the full-law
fiber precisely when
\(\underline\psi(P)=\overline\psi(P)\).  Conversely, if
\(\mathcal F(P)\ne\varnothing\), the construction (25) shows that \(P\) is
induced by a member of the class.  Thus (19) contains every and only observed
law whose full-law fiber identifies the target.  Any observable subclass with
that property is a subset of (19), which proves uniqueness and maximality.

The same elementary order argument shows
\[
 \underline\psi(P)<\overline\psi(P)
 \quad\Longleftrightarrow\quad
 \exists (\theta_1,x_1),(\theta_2,x_2)\in\mathcal F(P)
 \text{ with }T(\theta_1)\ne T(\theta_2).
\]
By (25)--(26), the two feasible points construct two compatible class members
with the same observed law and unequal targets.  Conversely, any such two
members give two such feasible points by the forward construction.  This is
the asserted exact two-SCM certificate.  It is finite: simplex membership,
the response sums, (14)--(16), and the two unequal objective values comprise
finitely many polynomial equalities, weak inequalities, strict inequalities,
and finite disjunctions.  The entries of \(B(P)\) and \(b(P)\) are rational
functions of observed cells whose denominators are positive by (13);
multiplication by those denominators turns their bridge equations into
polynomial equalities.  Thus an exhibited pair is checked directly without
any endpoint-attainment premise.

If \(\ell(P)\in\operatorname{row}(B(P))\),
\(\mathrm{thm{:}wbc\text{-}target\text{-}span\text{-}identification}\)
shows that every member of \(\mathcal C(P)\) has the same target.  If instead
the proxy-measurable arm-zero premise holds,
\(\mathrm{thm{:}proxy\text{-}measurable\text{-}arm\text{-}zero\text{-}rescue}\)
gives the same conclusion.  The endpoint characterization just proved puts
both sufficient subclasses inside \(\mathcal I_{\max}\).

By \(\mathrm{prop{:}scm\text{-}counterpair}\), \(P^{\dagger}\) is strictly
positive, its loading is not in the row space, and its full-law fiber contains
target values \(19/48\) and \(7/16\).  Exactness of (17) therefore gives
\[
 \overline\psi(P^{\dagger})-\underline\psi(P^{\dagger})
 \ge \frac7{16}-\frac{19}{48}=\frac1{24},
\]
which proves (20).

Finally, consider the universal positive-interior converse: at every induced
law with all observed cells positive, failure of row span should imply two
compatible class members with unequal targets.  Whenever
\(\mathcal F(P)\ne\varnothing\), feasibility supplies an \(x\) with
\(B(P)x=b(P)\), so \(b(P)\in\operatorname{col}(B(P))\).  By
\(\mathrm{lem{:}null\text{-}annihilator}\), row-span failure is then
equivalent to the existence of \(v\) with
\(B(P)v=0\) and \(\ell(P)^{\top}v\ne0\).  By (17), absence of two unequal
compatible targets is equivalent to constancy of \(T\) on
\(\mathcal F(P)\).  Consequently a counterexample to the universal converse
is exactly an element of \(\mathcal E_{\mathrm{int}}\).  The converse is
therefore true if and only if \(\mathcal E_{\mathrm{int}}=\varnothing\).
Nothing in this equivalence proves that emptiness; it isolates it as the
remaining latent-tilt question, exactly as claimed.
